\subsection{Plant model}

The motor is modelled by the first-order armature-to-speed relation used in
the reader. Electrical dynamics are neglected, so that the inertia \(J\),
viscous friction coefficient \(b\), motor constant \(K_m\), input \(u\), and
angular speed \(\omega\) satisfy
\begin{equation}
  J\dot{\omega}=-b\omega+K_m u .
  \label{eq:motor-differential-equation}
\end{equation}
The corresponding continuous-time transfer function is
\begin{equation}
  G(s)=\frac{\omega(s)}{u(s)}
      =\frac{K_m/J}{s+b/J}.
  \label{eq:motor-continuous-transfer}
\end{equation}

With zero-order hold sampling and sampling period \(T=\SI{0.08}{s}\), the
discrete transfer function is written in the notation of the reader as
\begin{equation}
  G(z)=\frac{b_0}{z+a_0},
  \label{eq:motor-discrete-transfer}
\end{equation}
where
\begin{equation}
  a_0=-\exp\!\left(-\frac{bT}{J}\right),
  \qquad
  b_0=\frac{K_m}{b}
      \left(1-\exp\!\left(-\frac{bT}{J}\right)\right).
  \label{eq:motor-discrete-parameters}
\end{equation}
The identified parameter vector is therefore
\begin{equation}
  x_e=\begin{bmatrix}a_0 & b_0\end{bmatrix}^{T}.
  \label{eq:identified-parameter-vector}
\end{equation}
